% Generated draft: requires mathematical and source review.
\documentclass[11pt]{article}
\usepackage[margin=25mm]{geometry}
\usepackage{fontspec}
\setmainfont{Latin Modern Roman}
\usepackage{amsmath,amssymb,mathtools,unicode-math}
\setmathfont{Latin Modern Math}
\usepackage{enumitem,needspace,xurl,hyperref}
\setlength{\parindent}{0pt}
\setlength{\parskip}{5pt}
\setlength{\emergencystretch}{3em}
\newcommand{\N}{\mathbb{N}}
\newcommand{\Z}{\mathbb{Z}}
\newcommand{\Q}{\mathbb{Q}}
\newcommand{\R}{\mathbb{R}}
\newcommand{\C}{\mathbb{C}}
\newcommand{\F}{\mathbb{F}}
\newcommand{\A}{\mathbb{A}}
\newcommand{\PP}{\mathbb{P}}
\newcommand{\E}{\mathbb{E}}
\newcommand{\eps}{\varepsilon}
\newcommand{\dd}{\,\mathrm{d}}
\begin{document}
\section*{UnsolvedMath ID: 30002266}
\textbf{Source problem code:} catalogue entry
\subsection{Definitions and mathematical statement}
A topological graph is understood here as a finite graph drawn in the plane: distinct vertices are represented by distinct points, and each edge is represented by a simple arc joining its two distinct endpoints, with its interior avoiding every vertex. Under the usual simplicity convention, two distinct edge arcs have at most one common point, which is either a common endpoint or a proper crossing of their interiors. In particular, loops and parallel edges are excluded. Two edges are called disjoint if their entire drawn images are disjoint; thus they share neither a vertex nor a crossing. These drawing and disjointness conventions are the conventional interpretation used in this draft; the supplied record does not explicitly define them, and their agreement with the cited report requires review.

\textbf{Question.} Does there exist an absolute constant \(C>0\) such that, for every pair of integers \(n\geq 1\) and \(k\geq 1\), every simple topological graph \(G\) with \(|V(G)|=n\) and no set of \(k\) pairwise disjoint edges satisfies
\[
|E(G)|\leq Ckn?
\]
Here the forbidden configuration consists of distinct edges \(e_1,\ldots,e_k\) whose drawn images satisfy
\[
e_i\cap e_j=\varnothing\qquad(1\leq i<j\leq k).
\]
The constant \(C\) is required to be independent of \(n\), \(k\), and the drawing. This is the uniform interpretation of the supplied \(O(kn)\) bound; the record does not separately state its quantifiers. When \(k=1\), the hypothesis forces \(E(G)=\varnothing\). The case \(n=0\), if allowed, is also trivial.

The supplied record labels the problem partially solved. Its dated assessment says that the requested bound remains open for simple topological graphs, that primary work gives only polylogarithmic upper bounds, and that linear bounds hold in geometric and special subclasses. The record provides no precise quantitative versions or subclass hypotheses for these auxiliary claims. They are supplied curation, not independently verified mathematical conclusions of this draft.
\subsection{Short English statement}
Is there a universal constant \(C\) such that every simple topological graph on \(n\) vertices with no \(k\) mutually disjoint drawn edges has at most \(Ckn\) edges? Here disjoint edges have neither shared endpoints nor intersections in the drawing. The drawing conventions and uniform interpretation of the bound require source review; this is a statement-only draft.
\subsection{Sources}
\begin{itemize}
\item Supplied local UnsolvedMath registry record, permanent ID \(30002266\), title: Disjoint-Edge Bounds for Simple Topological Graphs. Exact source problem code: \texttt{catalogue entry}. Collection: Oberwolfach Reports Open Problems. The record gives the proposed year as 2013 and names no individual proposer.
\item Supplied source citation: \emph{Graph Theory} (2013), Oberwolfach Reports, pp.\ 67--128, DOI \texttt{10.4171/owr/2013/02}. \url{https://doi.org/10.4171/owr/2013/02}. The supplied background describes Oberwolfach Reports as documenting mathematical workshops and their open-problem discussions.
\item Supplied literature source: \emph{Disjoint Edges in Topological Graphs}. \url{https://www.cs.bme.hu/~geza/disjoint.pdf}. No author names or further bibliographic details are supplied for this reference.
\item Supplied literature source: \emph{Graph Theory, Oberwolfach Report 02/2013}. \url{https://doi.org/10.4171/owr/2013/02}.
\item Literature assessment dated exactly \texttt{2026-08-21}: the supplied status is partially solved. The supplied assessment states that the desired \(O(kn)\) bound remains open for simple topological graphs, that primary work gives only polylogarithmic upper bounds, and that linear bounds hold in geometric and special subclasses. This assessment is supplied dated curation derived from the OWR curation, not independently verified here; time-sensitive status and cited claims require independent verification.
\item The registry supplies the statement label \texttt{corrected\_verified} and reports that the mathematical content was checked against the official Oberwolfach report, with notation, context, and question form normalized without changing the source-level claim. This is a supplied source-verification assertion; this draft does not claim to have visited or checked either reference.
\end{itemize}
\end{document}
